\section{The Three Interpretations Compared}
\label{sec:comparison}

The three interpretations answer three different questions raised by the same
sentence. The first asks what is owed to God and answers with the self, which
bears his image. The second asks how the debt to Caesar stands to God's
sovereignty and answers that it is real and bounded. The third asks what is
wrong with the questioners and answers that their hearts and words are
divided. The table sets out where each turns.

\begin{comparisontable}{}{The coin and the image}{Cyrus and Caesar}{The sincere heart}
Question answered & What are ``the things that are God's''? & How does the
debt to Caesar stand to God's sovereignty, and where does it end? & Why are
men who speak the truth called hypocrites, and what does Christ answer?\\
Decisive words & ``Whose image and inscription is this?'' (Mt 22:20) &
``Render \ldots\ to Caesar \ldots\ and to God'' (22:21), with Is 45:1, 4--6 &
``Why do you tempt me, ye hypocrites?'' (22:18), with 22:16\\
Principal inference & As the coin bears Caesar's image, man bears God's
(Augustine); Christ bears no image of Caesar, being God's image (Ambrose, on
Luke); the self is owed to its origin (Hilary) & God uses a gentile king who
does not know him (Jerome, Aquinas on Is 45); tribute is rendered as owed and
ends where godliness is harmed or Christ's cause is at stake (Chrysostom,
Augustine); joined under the Lectionary's pairing, Caesar's due is real and
bounded & The questioners' true words come from a divided heart (Jerome,
Chrysostom, Aquinas, Augustine); Christ knows and answers the heart as only God
does (Hilary, Aquinas), which is proof of his Godhead (Chrysostom)\\
Consequence & Self-gift, not a second tax & Civic duty within, and under, the
worship of God & The sincerity of the one who renders\\
Role of the first reading & Weakest: Jerome's reading of the Vulgate's
likeness, which the Lectionary's English lacks & Strongest: the official pairing, Cyrus as the Lord's
instrument & By contrast only: not knowing God and serving him, against
saying ``we know'' and tempting\\
Disagreement among the witnesses & Jerome's literal dues against the image & Hilary and
Ambrose (possession) against Chrysostom and Augustine (justice); the ``Christ
my Lord'' reading & Whether the reading of hearts proves Godhead (Chrysostom,
Hilary, Aquinas) or is any answerer's discernment (Jerome)\\
\end{comparisontable}

\subsection{What they share}

All three read the saying as setting Caesar's claim beneath God's, and none
reads it as a division of life into two equal realms, one political and one
religious. In the first, Caesar's claim reaches the coin and God's reaches the
one who carries it; in the second, Caesar's claim is real because God rules
the world in which it is made, and it ends where God's begins; in the third,
the debt to God includes the heart of the one who pays. All three rest on
Matthew 22:15--21 read as a whole scene, from the plot to the verdict. And all
three use the Collect's petition, for a will devoted to God and service of his
majesty with a sincere heart, but each for a different reason: the will as the
self that is owed (the first), the majesty served first (the second), the heart
undivided (the third).

All three also meet the counter-question in the same way. Augustine,
Chrysostom, Jerome and Aquinas agree that Jesus' question about the image is
not an evasion but the means of his answer, which he draws from the
questioners' own mouths. That agreement is developed in the third
interpretation, where it serves the argument most directly.

\subsection{Where they differ}

The interpretations differ in what they ask of a hearer. The first would be
satisfied by one who pays every tax and gives God himself. The second asks, in
addition, where the tax must be refused, and finds the answer in a soldier's
refusal to offer incense. The third asks whether the giver means what he says,
and would find wanting even a correct answer given with a divided heart.

They also differ in what they need from the rest of the Mass. Only the second
needs the first reading as an argument, and it alone stands on a relation that
the liturgical books declare. The first uses Isaiah only through Jerome's
reading of the Vulgate, and the third only by contrast. Conversely, the first
and the third depend most on the Missal's prayers, whose words about the will
and the heart they take up, while the second uses the prayers chiefly to place
Caesar's coin within an order of goods.

The first and the third can be preached together without strain: render
yourself, with an undivided heart. The second adds a public dimension that the
other two leave implicit, and with it a disagreement among the Fathers, about
why Caesar is owed anything at all, that the other two do not have to settle.
Each has its own principal difficulty: for the first, that the image of God is
an inference the saying does not state; for the second, that the Fathers speak to each
of the two readings the Lectionary pairs and not to the pair; for the third, that one of its own
witnesses, Jerome, treats the reading of hearts as a rule for any respondent
rather than a sign of Godhead.
